% Standalone problem document, generated from the adjacent README.md.
% Compile with XeLaTeX. Mathematical content is maintained in Markdown.
\documentclass[11pt,a4paper]{article}
\usepackage[fontsize=10.5pt]{scrextend}
\usepackage[margin=22mm,headheight=15pt,headsep=9mm,footskip=11mm]{geometry}
\usepackage{fontspec}
\setmainfont{texgyrepagella}[Extension=.otf,UprightFont=*-regular,BoldFont=*-bold,ItalicFont=*-italic,BoldItalicFont=*-bolditalic]
\setsansfont{texgyreheros}[Extension=.otf,UprightFont=*-regular,BoldFont=*-bold,ItalicFont=*-italic,BoldItalicFont=*-bolditalic]
\setmonofont{texgyrecursor}[Extension=.otf,UprightFont=*-regular,BoldFont=*-bold,ItalicFont=*-italic,BoldItalicFont=*-bolditalic,Scale=MatchLowercase]
\usepackage{amsmath,amssymb}
\usepackage{unicode-math}
\setmathfont{texgyrepagella-math.otf}
\usepackage{xcolor,microtype,enumitem,needspace,fancyhdr}
\usepackage{bookmark}
\definecolor{ink}{HTML}{17344B}
\definecolor{accent}{HTML}{197D80}
\definecolor{muted}{HTML}{596773}
\hypersetup{colorlinks=true,linkcolor=accent,urlcolor=accent,pdftitle={RA-03 - Improve
the randomized LU squared-error factor to
2\^{}k},pdfauthor={Open Problems in Numerical Linear Algebra}}
\urlstyle{same}
\setlength{\parindent}{0pt}
\setlength{\parskip}{5pt plus 1pt minus 1pt}
\linespread{1.04}
\setlength{\emergencystretch}{3em}
\widowpenalty=10000
\clubpenalty=10000
\displaywidowpenalty=10000
\allowdisplaybreaks[1]
\setlist{nosep,leftmargin=1.5em}
\providecommand{\tightlist}{\setlength{\itemsep}{0pt}\setlength{\parskip}{0pt}}
\providecommand{\pandocbounded}[1]{#1}
\pagestyle{fancy}
\fancyhf{}
\fancyhead[L]{\sffamily\footnotesize\color{muted}OPEN PROBLEMS IN NUMERICAL LINEAR ALGEBRA}
\fancyhead[R]{\sffamily\bfseries\footnotesize\color{accent}RA-03}
\fancyfoot[L]{\parbox[b]{0.9\textwidth}{\sffamily\fontsize{8}{10}\selectfont\color{muted}Editorial ratings; a bounded search cannot certify that a problem remains unsolved.\\Literature check: 2026-09-08}}
\fancyfoot[R]{\sffamily\footnotesize\color{muted}\thepage}
\renewcommand{\headrulewidth}{0.3pt}
\renewcommand{\headrule}{\hbox to\headwidth{\color{accent}\leaders\hrule height \headrulewidth\hfill}}
\makeatletter
\renewcommand{\section}{\Needspace{5\baselineskip}\@startsection{section}{1}{0pt}{12pt plus 2pt minus 2pt}{4pt}{\sffamily\bfseries\large\color{ink}}}
\renewcommand{\subsection}{\Needspace{4\baselineskip}\@startsection{subsection}{2}{0pt}{9pt plus 1pt minus 1pt}{3pt}{\sffamily\bfseries\normalsize\color{ink}}}
\makeatother
\setcounter{secnumdepth}{0}
\begin{document}
{\sffamily\small\color{muted}Randomized and low-rank approximation\par}
\vspace{3pt}
{\raggedright\hyphenpenalty=10000\exhyphenpenalty=10000\sffamily\bfseries\fontsize{21}{24}\selectfont\color{ink}Improve
the randomized LU squared-error factor to \(2^k\)\par}
\vspace{7pt}
{\sffamily\small\textbf{Difficulty:} challenging\quad\textbf{Importance:} interesting
to specialist\par}
\vspace{4pt}
{\color{accent}\hrule height 0.8pt}
\vspace{6pt}
\textbf{Topic:} randomized LU; low-rank approximation

\textbf{Status:} open; admitted after a bounded literature search found
no resolution.

\subsection{Problem statement}\label{problem-statement}

For \(A\in\mathbb C^{m\times n}\), define exact-arithmetic residuals
\(S_0=A\). At step \(t\), choose \((i,j)\) with conditional probability
\(|(S_t)_{ij}|^2/\|S_t\|_F^2\) and update

\[
S_{t+1}=S_t-\frac{S_t(:,j)S_t(i,:)}{(S_t)_{ij}}.
\]

After a zero residual, keep subsequent residuals zero. Prove or refute,
for every \(m,n\geq1\), every \(A\), and every \(1\leq k\leq\min(m,n)\),

\[
\mathbb E\|S_k\|_F^2\leq2^k
\sum_{j>k}\sigma_j(A)^2,
\]

where singular values decrease with \(j\). The bound concerns the mean
squared Frobenius error of this specified algorithm.

\subsection{References}\label{references}

Gilles and Wilber,
\href{https://arxiv.org/html/2601.22344v1}{\emph{Low-Rank Approximation
by Randomly Pivoted LU}}, Algorithm 1 and §3.1, Theorem 3, Eq. (10), and
the following conjecture (pp.~7--8).

\subsection{Status check}\label{status-check}

Their theorem gives \(4^k\); the conjecture explicitly replaces it by
\(2^k\). Searches for the title, \textit{randomly pivoted LU 2\^{}k},
and improved RPLU error bounds found no later resolution. The arXiv
submission history listed only v1 of January 29, 2026. The August 2026
RPCholesky theorem addresses a different pivot distribution and norm.
\end{document}
